% Estilo compartilhado: curriculo minimalista, duas colunas, sem foto.
% Incluido pelos arquivos de conteudo em pt/ e en/ via % Estilo compartilhado: curriculo minimalista, duas colunas, sem foto.
% Incluido pelos arquivos de conteudo em pt/ e en/ via % Estilo compartilhado: curriculo minimalista, duas colunas, sem foto.
% Incluido pelos arquivos de conteudo em pt/ e en/ via % Estilo compartilhado: curriculo minimalista, duas colunas, sem foto.
% Incluido pelos arquivos de conteudo em pt/ e en/ via \input{../common/style.tex}

\usepackage[a4paper,margin=1.4cm]{geometry}
\usepackage[T1]{fontenc}
\usepackage[utf8]{inputenc}
\usepackage{lmodern}
\usepackage{array}
\usepackage{enumitem}
\usepackage{xcolor}
\usepackage{hyperref}
\usepackage{paracol}

\definecolor{textgray}{gray}{0.35}

\pagestyle{empty}
\setlength{\parindent}{0pt}
\raggedright

\hypersetup{colorlinks=true, urlcolor=black, linkcolor=black}

% Colunas: barra lateral (esquerda) + conteudo principal (direita).
% paracol permite que cada coluna flua e quebre de pagina de forma
% independente, o que uma tabular normal nao suporta.
\columnratio{0.32}
\setlength{\columnsep}{22pt}

% Titulo de secao: versalete, com regra fina embaixo.
% Comando manual (nao usa \section/titlesec) porque \@startsection
% conflita com o controle de colunas do paracol.
\newcommand{\cvsection}[1]{%
  \par\vspace{8pt}%
  {\normalsize\scshape\bfseries #1}\par\vspace{2pt}%
  \noindent\rule{\linewidth}{0.4pt}\par\vspace{4pt}%
}

% Listas compactas para experiencias
\setlist[itemize]{leftmargin=*, itemsep=1pt, topsep=2pt, parsep=0pt}

% Nome + cargo no topo da barra lateral
\newcommand{\resumename}[2]{%
  {\Large\bfseries #1}\\[2pt]
  {\small\itshape\color{textgray}#2}\par\vspace{8pt}
}

% Item de experiencia: cargo, empresa, periodo alinhado a direita
\newcommand{\jobentry}[3]{%
  \textbf{#1}\hfill{\small\color{textgray}#3}\\
  {\small\itshape #2}\par\vspace{2pt}
}


\usepackage[a4paper,margin=1.4cm]{geometry}
\usepackage[T1]{fontenc}
\usepackage[utf8]{inputenc}
\usepackage{lmodern}
\usepackage{array}
\usepackage{enumitem}
\usepackage{xcolor}
\usepackage{hyperref}
\usepackage{paracol}

\definecolor{textgray}{gray}{0.35}

\pagestyle{empty}
\setlength{\parindent}{0pt}
\raggedright

\hypersetup{colorlinks=true, urlcolor=black, linkcolor=black}

% Colunas: barra lateral (esquerda) + conteudo principal (direita).
% paracol permite que cada coluna flua e quebre de pagina de forma
% independente, o que uma tabular normal nao suporta.
\columnratio{0.32}
\setlength{\columnsep}{22pt}

% Titulo de secao: versalete, com regra fina embaixo.
% Comando manual (nao usa \section/titlesec) porque \@startsection
% conflita com o controle de colunas do paracol.
\newcommand{\cvsection}[1]{%
  \par\vspace{8pt}%
  {\normalsize\scshape\bfseries #1}\par\vspace{2pt}%
  \noindent\rule{\linewidth}{0.4pt}\par\vspace{4pt}%
}

% Listas compactas para experiencias
\setlist[itemize]{leftmargin=*, itemsep=1pt, topsep=2pt, parsep=0pt}

% Nome + cargo no topo da barra lateral
\newcommand{\resumename}[2]{%
  {\Large\bfseries #1}\\[2pt]
  {\small\itshape\color{textgray}#2}\par\vspace{8pt}
}

% Item de experiencia: cargo, empresa, periodo alinhado a direita
\newcommand{\jobentry}[3]{%
  \textbf{#1}\hfill{\small\color{textgray}#3}\\
  {\small\itshape #2}\par\vspace{2pt}
}


\usepackage[a4paper,margin=1.4cm]{geometry}
\usepackage[T1]{fontenc}
\usepackage[utf8]{inputenc}
\usepackage{lmodern}
\usepackage{array}
\usepackage{enumitem}
\usepackage{xcolor}
\usepackage{hyperref}
\usepackage{paracol}

\definecolor{textgray}{gray}{0.35}

\pagestyle{empty}
\setlength{\parindent}{0pt}
\raggedright

\hypersetup{colorlinks=true, urlcolor=black, linkcolor=black}

% Colunas: barra lateral (esquerda) + conteudo principal (direita).
% paracol permite que cada coluna flua e quebre de pagina de forma
% independente, o que uma tabular normal nao suporta.
\columnratio{0.32}
\setlength{\columnsep}{22pt}

% Titulo de secao: versalete, com regra fina embaixo.
% Comando manual (nao usa \section/titlesec) porque \@startsection
% conflita com o controle de colunas do paracol.
\newcommand{\cvsection}[1]{%
  \par\vspace{8pt}%
  {\normalsize\scshape\bfseries #1}\par\vspace{2pt}%
  \noindent\rule{\linewidth}{0.4pt}\par\vspace{4pt}%
}

% Listas compactas para experiencias
\setlist[itemize]{leftmargin=*, itemsep=1pt, topsep=2pt, parsep=0pt}

% Nome + cargo no topo da barra lateral
\newcommand{\resumename}[2]{%
  {\Large\bfseries #1}\\[2pt]
  {\small\itshape\color{textgray}#2}\par\vspace{8pt}
}

% Item de experiencia: cargo, empresa, periodo alinhado a direita
\newcommand{\jobentry}[3]{%
  \textbf{#1}\hfill{\small\color{textgray}#3}\\
  {\small\itshape #2}\par\vspace{2pt}
}


\usepackage[a4paper,margin=1.4cm]{geometry}
\usepackage[T1]{fontenc}
\usepackage[utf8]{inputenc}
\usepackage{lmodern}
\usepackage{array}
\usepackage{enumitem}
\usepackage{xcolor}
\usepackage{hyperref}
\usepackage{paracol}

\definecolor{textgray}{gray}{0.35}

\pagestyle{empty}
\setlength{\parindent}{0pt}
\raggedright

\hypersetup{colorlinks=true, urlcolor=black, linkcolor=black}

% Colunas: barra lateral (esquerda) + conteudo principal (direita).
% paracol permite que cada coluna flua e quebre de pagina de forma
% independente, o que uma tabular normal nao suporta.
\columnratio{0.32}
\setlength{\columnsep}{22pt}

% Titulo de secao: versalete, com regra fina embaixo.
% Comando manual (nao usa \section/titlesec) porque \@startsection
% conflita com o controle de colunas do paracol.
\newcommand{\cvsection}[1]{%
  \par\vspace{8pt}%
  {\normalsize\scshape\bfseries #1}\par\vspace{2pt}%
  \noindent\rule{\linewidth}{0.4pt}\par\vspace{4pt}%
}

% Listas compactas para experiencias
\setlist[itemize]{leftmargin=*, itemsep=1pt, topsep=2pt, parsep=0pt}

% Nome + cargo no topo da barra lateral
\newcommand{\resumename}[2]{%
  {\Large\bfseries #1}\\[2pt]
  {\small\itshape\color{textgray}#2}\par\vspace{8pt}
}

% Item de experiencia: cargo, empresa, periodo alinhado a direita
\newcommand{\jobentry}[3]{%
  \textbf{#1}\hfill{\small\color{textgray}#3}\\
  {\small\itshape #2}\par\vspace{2pt}
}
